\input{lib/dissertationexclude.tex}

% Question: can spontaneous biotic selection effects stabilize substantially elevated functional complexity in an evolving system?

% In contrast to classical work concerning hypothetical plausibility of isolated scenarios, attempting to reproduce ETIs under controlled conditions tests the sufficiency of an integrated model to account for the origins of biological structure and complexity in concrete instances.
% Our work also contributes to broader debates in evolutionary biology over niche construction, which has emerged as a major tentpole of efforts to expand the modern synthesis \citep{DOTO}.

\section{Introduction}

Why should complex multicellular organisms exist, when simpler unicellular organisms grow biomass with substantially greater efficiency \citep{lynch2024bioenergetic}?
Adaptations in somatic division of labor and economies of scale \unskip\citep{Saunders1976,corning2015synergistic,niklas2014evolutionary,Bell1997,willensdorfer2009evolution,koschwanez2013improved} alone cannot fully explain complex multicellularity, as they recoup only a portion of metabolic costs; similarly, for explanations of somatic shielding \citep{gerhart1997cells,dawkins1976selfish}.
Some argue complex multicellular organisms to largely reflect passive diffusion from a lower complexity bound \citep{Gould1988,gould1990wonderful,gould2011full,mcshea2010biology,Lehman2010,Miconi2008}.
Under this neutralist framing, modest multicell population sizes act as a ratchet to accumulate physiological interdependencies that become locked into essential functionality \citep{lynch2003origins,Smith1972,stoltzfus1999possibility,muozgmez2021constructive,Liard2020}.

% this is a different type of ratchet, sentence bridge missing here: constructive neutral evolution at an ecosystem scale? mcshay zero force
Beyond constraints of metabolic viability, however, much of multicell adaptation tracks biotic interaction with other organisms --- the subject of niche construction theory, which examines how feedback from evolved organisms can shape selection \citep{odlingsmee2003niche,odlingsmee1988niche,levins1985dialectical,lewontin1983gene}.
% , through abiotic ecosystem engineering or direct production of biogenic structure.
Existing evolution experiments show niche construction can create spontaneous diversification --- for instance, nutrient crossfeeding among microbes \citep{poltak2010ecological,traverse2012tangled,kirisits2005characterization,ellis2015character,rosenzweig1994microbial,treves1998repeated}.
However, in addition to promoting diversity \citep{laland2017niche,brathen2015niche}, it is also hypothesized that niche construction effects can act as a selective ratchet for complex traits \citep{Waddington2008,laland2014niche,creanza2016cultural}.
Under this framing, selection for complex traits arises from evolving organisms themselves.
Existing agent-based simulation experiments, by imposing explicit biotic selection rules, can indeed produce such effects \citep{taylor2004niche,zaman2014coevolution,TODO}.
% Where advantaged by new environmental context, complex traits may be maintained by selection also act as a ratchet.

In the case of early multicellularity, vast morphological possibility opens broad new dimensions for biotic interaction \citep{Carroll2001,bonner1998origins}.
For instance, body size has long been recognized as means to resist predation \citep{niklas2014evolutionary,stanley1973ecological,bonner1988evolution}; such effects arise in experiments that introduce known microbial predators to populations of naive prey \citep{boraas1998phagotrophy,yoshida2003rapid,herron2019denovo,pentz2015predator}.
% multicellularity can shape ecological context and be sensitive to it (yeast) \citep{kalambokidis2023multispecies}.
% Indeed, current theories of complex multicellularity draw heavily on niche construction \citep{DOTO}.
% In the case of complex multicellularity, substantial observational evidence supports the evolutionary significance of niche construction \citep{DOTO} --- backed by equally rigorous theoretical developments \citep{DOTO}.
% we need to test this, but how do we test this without systems that can express it? a model where this CAN occur and is capable of it
To date, however, the role of \textit{de novo} niche construction in promoting complexity during early multicell evolution remains to be tested experimentally \citep{goldsby2012task,Joachimczak2012,goldsby2014evolutionary,bozdag2023denovo,ratcliff2012experimental,herron2016origins,kalambokidis2023multispecies}.

% Chicken and the egg problem --- temporal component, does it need the ability to continue changing the environment?
% what can we learn about niche construction that we can't through theory?
% concrete substantive instance --- what is unique about that? (Ratcliff?) --- complexity, novelty, adaptation

Leveraging a digital model system to study synthetic instances of multicell evolution \textit{in silico}, we observe spontaneous niche construction effects and test their relationship to evolved complexity.
Model digital organisms act analogously to biological cells: replicating, mutating, competing, and thus evolving \unskip\citep{pennock2007models}.
Though artificial in nature, the tractability of such digital models can nonetheless illuminate principles general to biological life \citep{adami2000evolution,lenski2003evolutionary,zaman2014coevolution} --- and beyond (for instance, principles of natural evolution have guided computational algorithms seeking ``open-ended'' creativity \citep{dolson2019constructive,soros2014identifying,huizinga2018emergence,lehman2012beyond,nguyen2015,stanley2019open,packard2019overview,taylor2016open,taylor2019evolutionary}).

The next section describes our experimental system: how model digital organisms replicate, form multicells, and mediate multicell-multicell interactions.
Following this, we examine a case study lineage to explore functional complexity and morphological novelty in early multicell evolution with greater detail than previously possible.
Competition assays reveal the case study's key morphological innovation as fit when abundant --- in other words, when shaping its own environmental context \citep{matthews2014under}.
We next test for biotic selection from this constructed niche.
% unexpected that self-selection doesn't "freeze" evolution but that
We find adaptation under biotic selection, and evolutionary replay experiments reveal it acting to lock in evolved morphological patterning and genotype complexity.
Finally, we test how niche construction and multicell complexity interact more generally, finding strong association across evolutionary replicates.

% Further competitions indicate sustained feedback between selection and niche construction, with self-benefit sensitive to time-shift manipulation.
% To assess whether
% Lastly, we present replay experiments that isolate the causal role of niche construction in maintaining complexity.


% outcomes between replay experiments with high self-context and low self-context, to identify the role of this self-context effect in maintaining evolved complexity.


% Although higher levels of complexity arose early in several other runs, they quickly went away.
% This led us to investigate why complexity was maintained in the case study strain.

% We found that this strain established its own biotic context.
% In replay experiments, we found that disrupting self-context of the focal strain led to reduced complexity.
% That is, that this strain was selecting on itself to maintain complexity.

% Here, we report , can stabilize substantially elevated functional complexity among \textit{de novo} multicells.
% Important evidence to the viability of niche construction perspectives on evolution of complexity, and open new to studying such processes in detail under controlled conditions.

% In this work, we apply a freeform \textit{in silico} modeling approach to evolve multicellular groups \textit{de novo} where life history unfolds in a shared environment.
% However, while existing evolution of multicellularity experiments have been conducted both \textit{in vivo} with X and Y and \textit{in silico} with agent-based modeling approaches... these have been characterized by experimental designs imposing strong abiotic selection effects with lesser opportunities for rich biotic interactions among multicells.

% The challenge and promise of open-ended evolution has animated decades of inquiry and discussion within the artificial life community \citep{packard2019overview}.
% The difficulty of devising models that produce continuing open-ended evolution suggests profound philosophical or scientific blind spots in our understanding of the natural processes that gave rise to contemporary organisms and ecosystems.
% Already, pursuit of open-ended evolution has yielded paradigm-shifting insights.
% For example, novelty search demonstrated how processes promoting non-adaptive diversification can ultimately yield adaptive outcomes that were previously unattainable \citep{lehman2011abandoning}.
% Such work lends insight into fundamental questions in evolutionary biology, such as the relevance --- or irrelevance --- of natural selection with respect to increases in complexity \citep{lehman2012evolution, Lynch8597} and the origins of evolvability \citep{lehman2013evolvability,Kirschner8420}.
% Evolutionary algorithms devised in support of open-ended evolution models also promise to deliver tangible, broader impacts for society.
% Possibilities include the generative design of engineering solutions, consumer products, art, video games, and AI systems .

% The preceding decades have witnessed advances in defining, both quantitatively and philosophically, the concept of open-ended evolution \citep{lehman2012beyond,dolson2019modes,bedau1998classification}.
% Notable investigations have also been made into causal phenomena that promote open-ended dynamics, such as ecological dynamics, selection, and evolvability \citep{dolson2019constructive,soros2014identifying,huizinga2018emergence}.
% The concept of open-endedness is fundamentally characterized by intertwined generation of novelty, functional complexity, and adaptation \citep{taylor2016open}.
% The strength and nature of relationships among these three phenomena, however, are yet to be fully understood.
% Here, we aim to complement ongoing work to develop a firmer theoretical understanding of the relationship between novelty, complexity, and adaptation by exploring the evolution of these phenomena through a case study using the DISHTINY digital multicellularity framework \dissertationexclude{\citep{moreno2019toward}}.
% We apply a suite of qualitative and quantitative measures to assess how these qualities can change over evolutionary time and in relation to one another.
